\documentclass[11pt]{article}
\usepackage[a4paper,margin=27mm]{geometry}
\usepackage{amsmath,amssymb,amsthm,mathtools}
\usepackage{graphicx}
\usepackage{microtype}
\usepackage[hidelinks]{hyperref}
\usepackage{enumitem}

\newtheorem{theorem}{Theorem}
\newtheorem{lemma}{Lemma}
\newcommand{\Ba}{\operatorname{Ba}}
\newcommand{\E}{\mathbb E}
\newcommand{\R}{\mathbb R}
\newcommand{\tp}{\widehat\otimes_\pi}

\title{Pointwise covariance determines the projective Dunford moment\\
\large An affirmative answer to Problem 9.5 of arXiv:2004.00237}
\author{Candidate solution packet}
\date{August 11, 2026}

\begin{document}
\maketitle

\noindent\textbf{Status.}
Candidate full solution, likely valid pending expert review.

\section{Source question}

Janson identifies $D([0,1]^m)$ isometrically with $C(K)$, where
$K=\widehat I^m$ and $\widehat I$ is the split interval.  The source
paper's Theorem~9.3 says that, for $D$-measurable $D([0,1]^m)$-valued
random variables $X,Y$ satisfying
\[
 M_X:=\sup_{t\in[0,1]^m}\E|X(t)|^2<\infty,
 \qquad
 M_Y:=\sup_{t\in[0,1]^m}\E|Y(t)|^2<\infty,
\]
equality of the pointwise second moments at all pairs of split points is
equivalent to equality of the injective Dunford second moments.  Theorem~8.8
also gives existence of the corresponding projective Dunford moments.
Problem~9.5 asks whether those projective moments must be equal
\cite{Janson}.

\begin{figure}[ht]
  \centering
  \includegraphics[width=\textwidth]{figures/open_problem_crop.png}
  \caption{Problem 9.5 on source PDF page 31.}
\end{figure}

Equivalently, the question is whether
\begin{equation}\label{eq:pointwise}
 \E[X(s)X(t)]=\E[Y(s)Y(t)]\qquad(s,t\in K)
\end{equation}
forces
\begin{equation}\label{eq:bilinear}
 \E\gamma(X,X)=\E\gamma(Y,Y)
\end{equation}
for every bounded bilinear form $\gamma$ on $C(K)$.

\section{Proof intuition}

Grothendieck's theorem extends any fixed bilinear form on $C(K)$ to a
Hilbert space $H=L^2(K,\nu)$.  The apparent obstruction is that a
$D$-measurable random function need not be jointly measurable on
$\Omega\times K$, so it is not immediate that it is a genuine random
variable in $H$.

The key observation is measure-theoretic rather than topological.  For every
finite Baire measure on $K$, its completed measure algebra is countably
generated.  Indeed, in each split-interval coordinate, the only information
beyond the usual Borel structure of $[0,1]$ is a countable correction; modulo
a fixed measure, only the countably many positive singleton atoms matter.
Thus $L^2(K,\nu)$ is separable.

The special Baire-Fubini theorem proved in the source paper then shows that all
scalar $H$-coordinates of $X$ are measurable and that $X$ has finite
Hilbert second moment.  Separability upgrades weak measurability to strong
measurability.  Applying the same Fubini theorem to the random function
$(s,t)\mapsto X(s)X(t)$ turns \eqref{eq:pointwise} into equality of all
matrix entries of the Hilbert covariance operators.  Those are positive
trace-class operators, hence the Hilbert projective second moments agree.
The Grothendieck extension then yields \eqref{eq:bilinear}.

\section{A separability lemma}

We write $\rho:\widehat I\to I=[0,1]$ for the map identifying the two copies
of each split point.  Proposition~A.1 of \cite{Janson} states that
\begin{equation}\label{eq:split-baire}
 \Ba(\widehat I)=
 \{\rho^{-1}(B)\mathbin\triangle N:
       B\in\mathcal B(I),\ N\subseteq\widehat I\text{ countable}\}.
\end{equation}
The same paper recalls that Baire sigma-fields commute with finite products.

\begin{lemma}\label{lem:sep}
Let $r\ge1$, let $K_r=\widehat I^r$, and let $\nu$ be a finite
positive Baire measure on $K_r$.  Then $L^p(K_r,\nu)$ is separable for
every $1\le p<\infty$.
\end{lemma}

\begin{proof}
Put $\Sigma_0=\Ba(\widehat I)$, so that
$\Ba(K_r)=\Sigma_0^{\otimes r}$.  Let $\nu_j$ be the $j$-th marginal
of $\nu$, and let
\[
 A_j:=\{x\in\widehat I:\nu_j(\{x\})>0\}.
\]
The set $A_j$ is countable.  Let $\mathcal G_j$ be the sigma-field
generated by all $\rho^{-1}(B)$, with $B$ ranging over a fixed countable
generator of $\mathcal B(I)$, together with the singletons $\{x\}$ for
$x\in A_j$.  Thus $\mathcal G_j$ is countably generated.

If $E\in\Sigma_0$, write $E=\rho^{-1}(B)\triangle N$ as in
\eqref{eq:split-baire}.  The countable set $N\setminus A_j$ is
$\nu_j$-null, while $N\cap A_j\in\mathcal G_j$.  Hence $E$ agrees
$\nu_j$-almost everywhere with a member of $\mathcal G_j$.

Set $\mathcal G=\mathcal G_1\otimes\cdots\otimes\mathcal G_r$, again a
countably generated sigma-field.  Every coordinate cylinder from
$\Sigma_0$ agrees $\nu$-almost everywhere with a cylinder from the
corresponding $\mathcal G_j$, because the inverse image of a
$\nu_j$-null set under the coordinate projection is $\nu$-null.  The
class of sets in $\Sigma_0^{\otimes r}$ which agree $\nu$-almost
everywhere with a member of $\mathcal G$ is a sigma-field.  It contains all
coordinate cylinders, and therefore equals $\Sigma_0^{\otimes r}$.

Thus the completion of $\Ba(K_r)$ modulo $\nu$ is already the completion
of a countably generated sigma-field.  Simple functions from a countable
generating algebra are dense in $L^p(K_r,\nu)$, proving separability.
\end{proof}

\section{Main result}

\begin{theorem}\label{thm:main}
Let $m\ge1$, $K=\widehat I^m$, and
$E=D([0,1]^m)=C(K)$.  Let $X,Y$ be $D$-measurable $E$-valued random
variables satisfying $M_X,M_Y<\infty$.  If \eqref{eq:pointwise} holds for
all $s,t\in K$, then \eqref{eq:bilinear} holds for every bounded bilinear
form $\gamma:E\times E\to\R$.  Consequently
\[
 \E[X\tp X]=\E[Y\tp Y]
 \quad\text{in }(E\tp E)^{**}.
\]
In particular, the answer to Problem~9.5 of \cite{Janson} is affirmative.
\end{theorem}

\begin{proof}
Fix a bounded bilinear form $\gamma$ on $E=C(K)$.  By the form of
Grothendieck's theorem used in the proof of \cite[Theorem~8.8]{Janson}, there
are a Baire probability measure $\nu$ on $K$ and a bounded bilinear form
$\widetilde\gamma$ on
\[
 H:=L^2(K,\nu)
\]
whose restriction to $C(K)\times C(K)$ is $\gamma$.

We first justify treating $X$ and $Y$ as square-integrable, strongly
measurable $H$-valued random variables.  Let $h\in H$, choose a Baire
representative, and define the finite signed Baire measure
$\lambda_h(A)=\int_A h\,d\nu$.  Since
$\sup_t\E|X(t)|\le M_X^{1/2}$, the Baire-Fubini theorem
\cite[Theorem~6.1(iii)]{Janson}, applied to $X$ and $\lambda_h$, shows
that
\[
 \omega\longmapsto \int_K X(\omega,s)h(s)\,d\nu(s)
       =\langle X(\omega),h\rangle_H
\]
is measurable.  Thus $X:\Omega\to H$ is weakly measurable.  Lemma
\ref{lem:sep} makes $H$ separable, so Pettis' measurability theorem makes
$X$ strongly measurable.  The same holds for $Y$.

Apply the positive part of the source's Baire-Fubini theorem to the
$D([0,1]^m)$-valued random variable $|X|^2$ and the measure $\nu$.
The source's extension bound from ordinary to split points gives
\[
 \E\|X\|_H^2
 =\int_K\E|X(s)|^2\,d\nu(s)
 \le M_X.
\]
Likewise $\E\|Y\|_H^2\le M_Y$.  Therefore the Bochner projective tensors
\[
 T_X:=\E[X\tp X],\qquad T_Y:=\E[Y\tp Y]
 \quad\text{in }H\tp H
\]
exist, since $\|x\otimes x\|_\pi=\|x\|_H^2$.

We next compare their matrix entries.  For $h,k\in H$, let
$\lambda_h=h\nu$ and $\lambda_k=k\nu$ as above.  The function
\[
 Z_X(\omega;s,t):=X(\omega,s)X(\omega,t)
\]
is a $D([0,1]^{2m})=C(K\times K)$-valued random variable: every sample is
continuous on $K\times K$, and each ordinary point evaluation is a product
of two measurable evaluations of $X$.  Moreover, Cauchy--Schwarz and the
split-point extension bound give
\[
 \sup_{s,t\in K}\E|Z_X(s,t)|\le M_X.
\]
Theorem~6.1(iii) of \cite{Janson}, now on $K\times K$ with the finite
signed Baire measure $\lambda_h\otimes\lambda_k$, yields
\begin{align*}
 \E[\langle X,h\rangle_H\langle X,k\rangle_H]
 &=\int_{K\times K}\E[X(s)X(t)]\,h(s)k(t)\,
       d\nu(s)d\nu(t).
\end{align*}
The same formula holds for $Y$.  Assumption \eqref{eq:pointwise} therefore
implies
\begin{equation}\label{eq:matrix}
 \E[\langle X,h\rangle_H\langle X,k\rangle_H]
 =\E[\langle Y,h\rangle_H\langle Y,k\rangle_H]
 \qquad(h,k\in H).
\end{equation}

Under the standard isometric identification of $H\tp H$ with the
trace-class operators on the real Hilbert space $H$, $T_X$ and $T_Y$
are the positive covariance operators whose matrix entries are the two sides
of \eqref{eq:matrix}.  Hence \eqref{eq:matrix} implies $T_X=T_Y$.
Finally, $\widetilde\gamma$ is a continuous functional on $H\tp H$, so
\[
 \E\gamma(X,X)
 =\langle\widetilde\gamma,T_X\rangle
 =\langle\widetilde\gamma,T_Y\rangle
 =\E\gamma(Y,Y).
\]
Since $\gamma$ was arbitrary, the projective Dunford moments in
$(E\tp E)^{**}$ are equal.  Their existence under the stated hypotheses is
also Theorem~8.8 of \cite{Janson}.
\end{proof}

\section{Consequences and limitations}

Combining Theorem~\ref{thm:main} with Theorem~9.3 of \cite{Janson}, all three
conditions are equivalent under the uniform pointwise second-moment bounds:
equality of split-point covariances, equality of injective Dunford second
moments, and equality of projective Dunford second moments.

The argument is specific to the second moments occurring in Problem~9.5.  It
does not prove that the canonical map
$(E\tp E)^{**}\to(E\widehat\otimes_\varepsilon E)^{**}$ is injective on
all of the bidual, the broader question in Remark~8.10.  It instead proves
injectivity on the particular class of differences of Dunford covariance
moments satisfying the source's pointwise $L^2$ bounds.

\section{Verification and novelty status}

The proof was checked specifically for four common hidden gaps: the
nonmetrizable Baire measure is reduced only modulo the actual Grothendieck
measure; weak measurability is upgraded only after $L^2$-separability is
proved; every double integration uses the source's special Fubini theorem on
$D([0,1]^{2m})$; and equality of elementary matrix entries is lifted through
the faithful trace-class representation of the Hilbert projective tensor.

On August 11, 2026, the local run indexes were searched by arXiv id, title,
and core terms.  Targeted web searches for the exact question and close
phrases found the source paper and the earlier background monograph
\cite{JansonKaijser}, but no later resolution.  The search was bounded, so
novelty confidence is moderate pending a specialist citation review.

\begin{thebibliography}{9}

\bibitem{Janson}
S. Janson,
\emph{The space $D$ in several variables: random variables and higher
moments},
Math. Scand. \textbf{127} (2021), no.~3, 544--584;
arXiv:2004.00237.

\bibitem{JansonKaijser}
S. Janson and S. Kaijser,
\emph{Higher moments of Banach space valued random variables},
Mem. Amer. Math. Soc. \textbf{238} (2015), no.~1127;
arXiv:1208.4272.

\bibitem{Grothendieck}
A. Grothendieck,
\emph{R\'esum\'e de la th\'eorie m\'etrique des produits tensoriels
topologiques},
Bol. Soc. Mat. S\~ao Paulo \textbf{8} (1953), 1--79.

\end{thebibliography}

\end{document}
